\documentclass[10pt,a4paper]{amsart}
\usepackage[margin=19mm]{geometry}
\usepackage{amsmath,amssymb,mathtools}
\usepackage[scr=rsfs,cal=euler]{mathalfa}
\usepackage{xfrac}
\usepackage[unicode,hidelinks,bookmarks=false]{hyperref}
\newcommand{\ie}{i.e.\ }
\newcommand{\cf}{cf.\ }
\newcommand{\resp}{respectively\ }
\newcommand{\Subsection}{\subsection}
\providecommand{\nobreakdash}{\nobreak\hbox{-}}
\setlength{\emergencystretch}{2em}
\multlinegap0pt
\usepackage{bm}
\usepackage{bbm}
\usepackage{dsfont}
\usepackage{xspace}
\usepackage{cancel}
\usepackage{mathtools}
\usepackage{amsthm}
\makeatletter
\@ifundefined{lemma}{\newtheorem{lemma}{Lemma}}{}
\@ifundefined{theorem}{\newtheorem{theorem}{Theorem}}{}
\makeatother
\title[A CAT(0)-approach to the marked length spectral rigidity of ]{A CAT(0)-approach to the marked length spectral rigidity of Sinai billiards}
\author{Douglas Finamore, Martin Leguil}
\date{Source: arXiv:2510.18983v1 (CC BY 4.0)}
\begin{document}
\maketitle

\noindent\emph{Source and licence.} A CAT(0)-approach to the marked length spectral rigidity of Sinai billiards. Version arXiv:2510.18983v1 (CC BY 4.0), distributed under the Creative Commons Attribution 4.0 International licence, as recorded in the arXiv licence metadata for this exact version and shown on the abstract page https://arxiv.org/abs/2510.18983v1. Full rights notice: licenses/arxiv-2510.18983-CC-BY-4.0.txt.\par\smallskip

\noindent\textit{Editorial scope.} Counted unit: the Theorem whose source label is \texttt{thm:char\_billiard\_cycles\_1} in the article (source0.tex lines 2015--2017, source label \texttt{thm:char\_billiard\_cycles\_1}), delivered together with the article's own complete original proof of it (source0.tex lines 2018--2076, one \texttt{proof} environment of 59 lines). No mathematics is written, repaired or re-derived: statement and proof are the article's own text line for line. Required context retained uncounted: thm:K\_0\_as\_gluing\_metric\_space (lines 650--656); lem:unique\_minimising\_geodesic (lines 1180--1182). Every definition, display and label of those lines is retained unchanged. Omitted from this package and not redistributed: 1--649, 657--1179, 1183--2014, 2077--2365. Nothing inside a retained excerpt is omitted or altered: every retained body is delivered exactly as the article's lines are written, so no omission receipt of any kind accompanies this package. Citations: the retained excerpts carry no citation command at all, so no bibliography entry of the article is transcribed and no reference is redistributed. Reference adaptation: none was needed and none was made; every reference inside the delivered unit resolves inside it, so no unresolved reference or citation is produced. Printed slips kept literal: no printed word, symbol, hypothesis or number of the counted statement or of its proof was altered. Layout and class: the article's own class is \texttt{amsart} with packages including presets; the wrapper is the lane's pinned \texttt{amsart} editorial wrapper at 10pt on A4, which restates the theorem-like environments the retained text needs and adds the article's own macro definitions that the retained text needs (none), with extra wrapper packages (bm, bbm, dsfont, xspace, cancel, mathtools). The delivered numbering is the unit's own. Assets: no retained span contains a figure environment, a graphics-inclusion command, a PDF-page inclusion, a table or a TikZ picture; no image file is copied into this package, the package declares no asset dependency and no image file is redistributed with it. Licence: version arXiv:2510.18983v1 of this article is distributed under the Creative Commons Attribution 4.0 International licence, as recorded in the arXiv licence metadata for this exact version and shown on the abstract page https://arxiv.org/abs/2510.18983v1. A full rights notice is delivered at licenses/arxiv-2510.18983-CC-BY-4.0.txt. Characters: every retained excerpt is pure ASCII, so the delivered engine, XeLaTeX with its default Latin Modern text font, reports no missing character.
\begin{theorem}\label{thm:K_0_as_gluing_metric_space}
    Suppose \( K^\circ \) has \( k \) distinct connected components. Then the space \( K_0 \) is isometric to the metric quotient 
    \[
        ^{\textstyle \left(\displaystyle{\coprod_{i=1}^k}(\mathcal{D}, d_\mathcal{L})\right)}\!\!\!\Big/_{\!\!\textstyle\sim \ }
    \]
    where \( \sim \) is the equivalence relation identifying corresponding boundary components.
\end{theorem}

\begin{lemma}\label{lem:unique_minimising_geodesic}
    Every element \( \alpha \in \mathcal{C}(K) \) contains a single closed geodesic curve \( \gamma^\alpha_0 \), which is a curve of minimal \(d_0\)-length among all those of \( \alpha \).
\end{lemma}

\begin{theorem}\label{thm:char_billiard_cycles_1}
    There is a surjection, induced by \( p \), between the set \( \mathcal{C}(K) \) and the collection of all billiard cycles.
\end{theorem}

\begin{proof}
    \par We begin by showing that, for any \( \alpha \in \mathcal{C}(K) \), the geodesic \( \gamma^\alpha_0 \) projects to a billiard cycle.
    We consider an arc-length parametrisation \( \gamma^\alpha_0\colon [0, l_\alpha] \to K_0 \).
    We need to construct a partition \( \{t_0 < \cdots < t_{k-1}\} \) of \( [0, l_\alpha] \) satisfying properties (i), (ii) and (iii).
    \par First, suppose \( \gamma^\alpha_0 \subset B(\gamma^\alpha_0) \).
    In this case, the curve \( \gamma^\alpha_0 \) coincides with one the obstacles \( \Gamma_i \), the trace of which is already a geodesic of \( d_\mathcal{L} \).
    We simply take \( t_0 = 0 \) and it follows that \( p(\gamma^\alpha_0) \) has the desired properties.
    \par Otherwise, the set \( \gamma^\alpha_0\setminus B(\gamma^\alpha_0) = K^\circ\cap\gamma^\alpha_0 \) is a non-empty open subspace of \( \gamma^\alpha_0\).
    Hence there is a finite collection of open subintervals \(I_j \subset [0, l_\alpha] \) such that the decomposition of  \( K^\circ\cap\gamma^\alpha_0 \) into connected components is
    \[
        K^\circ\cap\gamma^\alpha_0 = \coprod_{j=0}^{q}\gamma^\alpha_0(I_j)
    \]
    We assume our intervals are labelled in order, such that \( \sup I_j < \inf I_{j+1} \).
    Let \( t_j := \sup I_j \). 
    By construction \( \gamma^\alpha_0(t_j) \) is a collision point and, for any \( t \in I_j \), the triple \( (t_j-t; p(\gamma^\alpha_0(t))) \) is an element of \( A_0 \).
    It follows that each \( \gamma^\alpha_0\rvert_{I_j} \) projects to a straight line on \( \mathcal{D} \).
    Moreover, if the collision at \(t_j\) is non-degenerate, then \(\gamma^\alpha_0(t_j)\) is an isolated element of \(B(\gamma^\alpha_0)\), hence \(\inf I_{j+1} = \sup I_j = t_j\).
    \par Thus, at a non-degenerate collision \( \gamma^\alpha_0(t_j) \) the set
    \[
        p(\gamma^\alpha_0\rvert_{I_j})\cup p(\gamma^\alpha_0\rvert_{I_{j+1}})
    \]
    is the union of two geodesic segments in \( (\mathcal{D}, d_\mathcal{L}) \) with a common end point. If the collision at \(t_j\) is grazing, then \( p(\gamma^\alpha_0(t)\rvert_{I_j \cup I_{j+1}}) \) is actually a straight line. 
    On the other hand, if it is regular, then the geodesic \( \gamma^\alpha_0 \) is still conjugated to the geodesic flow after the collision, i.e., \(\inf I_{j+1} = t_j \) and \( (t_{j+1}-t; p(\gamma^\alpha_0(t))) \in A_0 \) for any \( t\in I_j \).
    In particular, \( \gamma^\alpha_0(t)\rvert_{I_j \cup I_{j+1}} \) projects to a billiard trajectory.
    In any case, we see that when the collision point at \(t_j\) is non-degenerate, the trajectory respect the desired specular reflection law at \( p(\gamma^\alpha_0(t)\rvert_{I_j \cup I_{j+1}}) \) at \(t_j\).
    \par If, on the other hand, the collision point at \(t_j\) is degenerate, then \(s_{j}:=\inf I_{j+1} > t_j \) and \(p(\gamma^\alpha_0\rvert_{[t_j, s_j]})\) is a segment of the boundary \( \partial \mathcal{D} \). 
    Now, we notice that in this case \(\gamma^\alpha_0(I_j \cup [t_j, s_j] \sup I_{j+1}\) is entirely contained in the closure of a connected component of \( K^\circ \), hence  \(p\) maps it into a geodesic of \( (\mathcal{D}, d_\mathcal{L}). \)   
    \par So we proceed as follows: the first element of our partition is \( t_0 \). 
    If \( t_1 \) is a regular collision point, we add it to the partition.
    If not, we look at the next regular collision point \( t_j \). 
    The segments between regular collisions points are unions of straight lines and boundary segments, and such unions are geodesics because they are images under \( p \) of segments of \( \gamma^\alpha_0 \). 
    As regular collision points obey the specular reflection law, the collection of all of them forms the desired partition.  
    This finishes the first part of the proof.
    \par Conversely, we show that every billiard cycle of \( \mathcal{D} \) is the image under \( p\) of a geodesic \( \gamma^\alpha_0 \).
    Let \( \sigma \) be a billiard cycle, and \( x_i = \sigma(t_i), \ i = 0, 1, \cdots, k-1, k \) its regular collision points.
    As before, we let \( t_k = t_0 \).
    We lift \( \sigma \) to a curve \( \sigma_\mathrm{up} \) in \( K_0 \) in the following way.
    For each \( x_i \in \partial\mathcal{D}\) there is a single point \( \widetilde x_i \in B(K) \) which projects to \( x_i \), since \( p\rvert_{B(K)} \) is a diffeomorphism.
    We lift the first geodesic segment \( \sigma\rvert_{[t_0, t_1]} \) to the sheet \( \mathcal{D}_\mathrm{up} \).
    This can be done in a unique way, since \( p\rvert_{\mathcal{D}_\mathrm{up}}: (\mathcal{D}_\mathrm{up}, d_0) \to (\mathcal{D}, d_\mathcal{L}) \) is an isometry.
    Now, in the same way, we lift the second geodesic segment \( \sigma\rvert_{[t_1, t_2]} \), this time to the sheet \( \mathcal{D}_\mathrm{down} \).
    As we showed in Theorem \ref{thm:K_0_as_gluing_metric_space}, if geodesic segments from different sheets are joined a collision point, then the curve obtained is minimising exactly when the collision is specular.
    Hence the lifting of \( \sigma\rvert_{[t_0, t_2]} \) thus defined is a geodesic segment of \( K_0 \).
    \par If we proceed this way, we get a closed curve \( \sigma_\mathrm{up} \) in \(K_0 \) which is everywhere minimising, except maybe at \( t_0 \).
    Indeed, if the lifts \( \sigma\rvert_{[t_0, t_1]} \) and \( \sigma\rvert_{[t_0, t_2]} \) are in the same sheet\footnote{
        this happens exactly when we have an even number \( k \) of regular collision times \( t_i \), or equivalently an odd number of collision points \( x_i \)
    },
    then \( \sigma\rvert_{[t_k, t_1]} \) is not minimising around \( t_0 \):
    convexity of the scatterers means that for any two points \( \sigma(t_0-\epsilon) \) and \( \sigma(t_0 + \epsilon) \) the line
    \[
        r(t) = \sigma(t_0-\epsilon) + t(\sigma(t_0+\epsilon) + \sigma(t_0-\epsilon)) 
    \]
    is entirely contained in the interior of \( \mathcal{D} \) and has Euclidean length smaller than the sum \( l_0(\sigma\rvert_{[t_0-\epsilon, t_0]}) + l_0(\sigma\rvert_{[t_0, t_0+\epsilon]}) \).
    To remedy this, we can consider the curve \( \sigma_\mathrm{down} \), which is the lift of \( \sigma \) to \( K_0 \) constructed as before, but starting at the sheet \( \mathcal{D}_\mathrm{down} \).
    Then the curve \( \sigma_\mathrm{up}\ast\sigma_\mathrm{down}^{-1} \), by construction, is a concatenation of geodesic segments, all of whose regular collisions are specular, and is therefore a geodesic of \( K_0 \).
    \par Finally, we let \( \alpha \) be the free homotopy class of \( \sigma_\mathrm{up}\ast\sigma_\mathrm{down}^{-1} \). 
    Since \( \sigma_\mathrm{up}\ast\sigma_\mathrm{down}^{-1} \) is a closed geodesic, it follows from Lemma \ref{lem:unique_minimising_geodesic} that \(\gamma^\alpha_0 = \sigma_\mathrm{up}\ast\sigma_\mathrm{down}^{-1} \).
    Thus \(p(\gamma^\alpha_0) = \sigma \), as we wanted.  
\end{proof}

\end{document}
